\part{Finite typed transfer and the formal contract}\label{part:finite}
\noindent The finite release supplies the algebraic foundation and a contract for its individual declarations. Here, ``roadmap'' and ``outside the contract'' describe the scope of the finite Lean formalization. Ordinary mathematical results in Part~II may lie beyond that scope. The thirty original result environments and their proof bodies are retained. The source contract has 41 mapping rows because it also includes definitions and other interfaces.

\medskip\noindent The preserved v1.0.1 release records Lean v4.34.0-rc2 against Mathlib revision \leanname{aeb72f1e9981b6d35b426b16dca7759be3e7ba06} \cite{lean,mathlib}. The compiler and kernel check descriptions in Part~I refer to that inherited evidence. The companion verification files record the editorial checks of this edition separately.
\section{Purpose and formal status}
\label{sec:purpose}

Paper I fixed typed seeds, operator grammars, exact projector packets,
generated hulls, and reconstruction \cite{paperI}.  Paper II added coherent
refinement coordinates, weighted transfer, arithmetic separation, and linear
defect identities \cite{paperII}.  The present paper begins the next layer:
explicit edge words, optional connection data, and obstructions within
individual packets.

Canonical coordinate transport and a connection play different mathematical
roles. Canonical transport compares two exact descriptions of the
\emph{same} ambient object and is forced by analysis and synthesis. A
connection chooses transports that may carry extra modeling, discretization,
or gauge information. Canonical transport obeys a cocycle theorem. A
connection may fail to obey a cocycle, and that failure is meaningful data.
Keeping separate notation lets us distinguish the automatic coherence of
exact coordinates from the possible holonomy of a chosen connection.

\subsection{Statement labels}

Every substantial theorem, lemma, corollary, proposition, definition, API,
and roadmap item carries one or more of the three labels below. When separate
clauses use a Paper III result and an inherited bridge, the statement carries
both applicable labels. Illustrative examples do not require a label.
\begin{enumerate}[label=\textbf{S\arabic*.},leftmargin=3.2em]
\item \textbf{Lean-backed in Paper III}: the displayed scope is no stronger
than the named compiled declaration in the released Paper III module.
\item \textbf{Inherited Lean-backed}: the displayed result is supplied by the
frozen formal contract of Paper I or II.
\item \textbf{Roadmap only}: the text states a future interface or theorem
target and is excluded from the completion denominator.
\end{enumerate}

\begin{boundary}[Final theorem denominator]
The completion denominator of the finite v1.0.1 contract consists of the
thirty result environments in Part~I. Of these, 26 resolve to compiled
Paper III declarations and four resolve to the frozen Papers I--II
contracts. The authoritative theorem map names the exact declarations used
by each statement. The kernel audit checks those targets and their supporting
theorems. The namespace gate requires every declaration in the compiled
\leanname{Komorebi.PaperIII} namespace to pass its checks. Roadmap items are
excluded from the denominator and remain visibly marked.
\end{boundary}

\section{Finite typed packet graphs}
\label{sec:edge-operators}

\subsection{Orientation and matrix convention}

Let $V$ be a finite packet type, $E$ a finite edge type, and $R$ a semiring.
A directed multigraph is a pair of endpoint functions
\[
  \src,\tgt:E\longrightarrow V.
\]
Parallel edges remain distinct elements of $E$.  Rows are targets and columns
are sources.  Thus an edge $e:i\to j$ is represented by a matrix supported at
row $j$, column $i$.

\begin{definition}[Packet graph and scalar edge system]
\label{def:edge-system}
For $a,b\in V$ and $r\in R$, define
\[
  \operatorname{packetMatrix}(a,b;r)_{yx}
  =
  \begin{cases}
  r,&y=b\text{ and }x=a,\\
  0,&\text{otherwise}.
  \end{cases}
\]
An edge system is a coefficient function $c:E\to R$.  Its elementary edge
operator and total transfer operator are
\[
  L_e=\operatorname{packetMatrix}(\src(e),\tgt(e);c(e)),
  \qquad
  T=\sum_{e\in E}L_e.
\]
\end{definition}
\statusleanthree{Komorebi.PaperIII.EdgeOperators.PacketGraph; Komorebi.PaperIII.EdgeOperators.PacketGraph.packetMatrix; Komorebi.PaperIII.EdgeOperators.PacketGraph.EdgeSystem; Komorebi.PaperIII.EdgeOperators.PacketGraph.EdgeSystem.edgeOp; Komorebi.PaperIII.EdgeOperators.PacketGraph.EdgeSystem.transferOp}

Words are written in execution order. In $L_fL_e$, the rightmost operator
$L_e$ acts first. For a written word $(e_1,\ldots,e_n)$, $e_1$ acts first
and the operator product is $L_{e_n}\cdots L_{e_1}$.

\begin{keyidea}[title={Reading the execution order}]
Read the edge list in the order of execution. The matrix product writes later
actions farther to the left. The target packet of one edge must equal the
source packet of the next before multiplication can contribute. Over a noncommutative
semiring, this order also fixes the order of the scalar coefficients.
\end{keyidea}

\begin{lemma}[Declared support]
\label{lem:edge-support}
For every edge $e$, $(L_e)_{yx}=0$ unless
$y=\tgt(e)$ and $x=\src(e)$.
\end{lemma}
\statusleanthree{Komorebi.PaperIII.EdgeOperators.PacketGraph.EdgeSystem.edgeOp_eq_zero_off_support}

\begin{theorem}[Typed interface annihilation]
\label{thm:edge-annihilation}
If $\tgt(e)\ne\src(f)$, then
\[
  L_fL_e=0.
\]
\end{theorem}
\statusleanthree{Komorebi.PaperIII.EdgeOperators.PacketGraph.EdgeSystem.edgeOp_mul_edgeOp_eq_zero}

\begin{proof}
For a matrix entry $(y,x)$,
\[
  (L_fL_e)_{yx}=\sum_{z\in V}(L_f)_{yz}(L_e)_{zx}.
\]
A potentially nonzero summand would require simultaneously
$z=\src(f)$ and $z=\tgt(e)$, contrary to the hypothesis.
\end{proof}

\begin{theorem}[Collapse of a matched edge pair]
\label{thm:edge-composable}
If $\tgt(e)=\src(f)$, then
\[
L_fL_e
=
\operatorname{packetMatrix}
\bigl(\src(e),\tgt(f);c(f)c(e)\bigr).
\]
Together with \cref{thm:edge-annihilation}, this gives both cases for
the composition of two edge operators.
\end{theorem}
\statusleanthree{Komorebi.PaperIII.EdgeOperators.PacketGraph.EdgeSystem.edgeOp_mul_edgeOp_of_composable}

\begin{proof}
The matrix-product sum has exactly one potentially nonzero intermediate
index, namely $\tgt(e)=\src(f)$.  At that index the coefficients occur in
operator order as $c(f)c(e)$, and the outer source and target are
$\src(e)$ and $\tgt(f)$.
\end{proof}

\begin{boundary}[Current fiber scope]
The compiled Paper III model above has one scalar coordinate per packet.
For a family of modules $(X_i)_{i\in V}$, an intrinsic edge map would instead
have type $\ell_e:X_{\src(e)}\to X_{\tgt(e)}$, and its totalization on
$\bigoplus_iX_i$ would be
$\iota_{\tgt(e)}\ell_e\pi_{\src(e)}$.  The same support proof applies, but
this generalization with dependent fibers is not attributed to the released
Paper III modules for scalar edge systems.
\end{boundary}

\Cref{fig:typed-edge-walks} summarizes the support convention and shows
how composable words differ from words with a broken interface. The next
section uses this distinction.

\begin{figure}[tbp]
\centering
\includegraphics[width=.96\linewidth]{typed_edge_walks.pdf}
\caption{Rows record target packets and columns record source packets.
This convention fixes execution order. Composable words collapse to endpoint
blocks, while one broken packet interface annihilates the entire operator
product.}
\label{fig:typed-edge-walks}
\end{figure}

\section{Operator words and directed walks}
\label{sec:words-walks}

\subsection{Typed words}

\begin{definition}[Composable word and execution operator]
\label{def:word-op}
A finite word $w=(e_1,\ldots,e_n)$ is composable when
\[
  \tgt(e_r)=\src(e_{r+1})
  \qquad (1\le r<n).
\]
Set
\[
  L_w=L_{e_n}\cdots L_{e_1},
  \qquad L_{\varnothing}=I.
\]
This definition records the written order of execution even though matrix
composition is read right to left.
\end{definition}
\statusleanthree{Komorebi.PaperIII.EdgeOperators.PacketGraph.EdgeSystem.Composable; Komorebi.PaperIII.EdgeOperators.PacketGraph.EdgeSystem.wordOp}

\begin{theorem}[First broken interface]
\label{thm:first-break}
If a word begins $(e,f,\ldots)$ and $\tgt(e)\ne\src(f)$, then its word
operator is zero.
\end{theorem}
\statusleanthree{Komorebi.PaperIII.EdgeOperators.PacketGraph.EdgeSystem.wordOp_eq_zero_of_first_break}

\begin{theorem}[Noncomposable words vanish]
\label{thm:nonwalk-vanish}
If $w$ is not composable, then $L_w=0$.
\end{theorem}
\statusleanthree{Komorebi.PaperIII.EdgeOperators.PacketGraph.EdgeSystem.wordOp_eq_zero_of_not_composable}

\begin{proof}
Choose the first broken adjacent pair.  The corresponding consecutive factor
is zero by \cref{thm:edge-annihilation}, and associativity absorbs that zero
into the complete product.  The Lean proof implements the same argument by
induction on the list.
\end{proof}

\subsection{The checked recursive expansion}

\begin{definition}[Recursive word expansion]
\label{def:recursive-word-expansion}
Define
\[
  \cW_0=I,
  \qquad
  \cW_{n+1}=T\cW_n.
\]
\end{definition}
\statusleanthree{Komorebi.PaperIII.EdgeOperators.PacketGraph.EdgeSystem.wordExpansion}

\begin{theorem}[Transfer power equals recursive word expansion]
\label{thm:power-recursive-expansion}
For every $n\in\N$,
\[
  T^n=\cW_n.
\]
\end{theorem}
\statusleanthree{Komorebi.PaperIII.EdgeOperators.PacketGraph.EdgeSystem.transferOp_pow_eq_wordExpansion}

\begin{proof}
The zero case is the identity matrix.  If $T^n=\cW_n$, then the
left-successor power convention gives
$T^{n+1}=TT^n=T\cW_n=\cW_{n+1}$.
\end{proof}

\begin{remark}
The formal module also defines an explicit sum over tuples
$w:\operatorname{Fin}(n)\to E$ and a second sum in which noncomposable tuples
are replaced by zero.  It proves that the recursive matrix, the explicit tuple sum, and the sum
filtered by composability agree. The separate declarations make the finite
enumeration claim directly checkable.
\end{remark}
\statusleanthree{Komorebi.PaperIII.EdgeOperators.PacketGraph.EdgeSystem.explicitWordExpansion; Komorebi.PaperIII.EdgeOperators.PacketGraph.EdgeSystem.walkExpansion; Komorebi.PaperIII.EdgeOperators.PacketGraph.EdgeSystem.wordExpansion_eq_explicitWordExpansion; Komorebi.PaperIII.EdgeOperators.PacketGraph.EdgeSystem.explicitWordExpansion_eq_walkExpansion}

\subsection{Exact expansion from words to walks}

For a nonempty composable word $w=(e_1,\ldots,e_n)$, define its reverse
coefficient product
\[
  c^\leftarrow(w)=c(e_n)c(e_{n-1})\cdots c(e_1).
\]
The reversal is essential over a noncommutative semiring.

\begin{lemma}[One word is one endpoint block]
\label{lem:word-endpoint}
If $w=(e_1,\ldots,e_n)$ is nonempty and composable, then
\[
  L_w
  =
  \operatorname{packetMatrix}
  \bigl(\src(e_1),\tgt(e_n);c^\leftarrow(w)\bigr).
\]
\end{lemma}
\statusleanthree{Komorebi.PaperIII.EdgeOperators.PacketGraph.EdgeSystem.reverseCoefficient; Komorebi.PaperIII.EdgeOperators.PacketGraph.EdgeSystem.terminalTarget; Komorebi.PaperIII.EdgeOperators.PacketGraph.EdgeSystem.wordOp_eq_packetMatrix_of_composable}

\begin{proof}
Induct on $n$.  The case $n=1$ is the definition of $L_e$.  Suppose the
formula holds for the prefix $(e_1,\ldots,e_n)$ and append $f$.  Composability
gives $\tgt(e_n)=\src(f)$.  Multiplying the two packet matrices leaves one
intermediate index and changes the coefficient to
$c(f)c(e_n)\cdots c(e_1)$, with source $\src(e_1)$ and target $\tgt(f)$.
\end{proof}

Write $\Walk_n(i,j)$ for the finite set of composable edge tuples of length $n$
whose first source is $i$ and whose final target is $j$. For $n=0$, it has
one empty walk when $i=j$ and is empty otherwise. Assign the empty walk the
coefficient $1$.

\begin{interpretation}[title={The empty walk}]
The empty walk records doing nothing at a chosen packet. It contributes only
to a diagonal entry. Thus the convention for length zero matches the identity
matrix and lets the walk sum include every nonnegative power.
\end{interpretation}

\begin{theorem}[Finite operator expansion from words to walks]
\label{thm:word-to-walk}
For every $n\in\N$,
\[
  \boxed{
  T^n
  =
  \sum_{w\in E^n}L_w
  =
  \sum_{w\in E^n}
  \begin{cases}
  L_w,&w\text{ composable},\\
  0,&w\text{ not composable}.
  \end{cases}
  }
\]
Parallel edges occur as distinct summands.
\end{theorem}
\statusleanthree{Komorebi.PaperIII.EdgeOperators.PacketGraph.EdgeSystem.transferOp_pow_eq_explicitWordExpansion; Komorebi.PaperIII.EdgeOperators.PacketGraph.EdgeSystem.transferOp_pow_eq_walkExpansion}

\begin{proof}
Induction and finite distributivity give
\[
\left(\sum_{e\in E}L_e\right)^n
=
\sum_{(e_1,\ldots,e_n)\in E^n}L_{e_n}\cdots L_{e_1}.
\]
Every noncomposable summand is zero by
\cref{thm:nonwalk-vanish}, so inserting the composability filter does not
change the sum.  In Lean, tuple concatenation and the successor sum are
handled by \leanname{tupleWordOp_snoc} and
\leanname{explicitWordExpansion_succ}; the two bridge equalities then yield
the displayed theorem.
\end{proof}

\begin{corollary}[Endpoint coefficient walk sum]
\label{cor:word-to-walk-entry}
For all $n\in\N$ and $i,j\in V$,
\[
  \boxed{
  (T^n)_{ji}
  =
  \sum_{w\in\Walk_n(i,j)}c^\leftarrow(w).
  }
\]
\end{corollary}
\statusleanthree{Komorebi.PaperIII.EdgeOperators.PacketGraph.EdgeSystem.Connects; Komorebi.PaperIII.EdgeOperators.PacketGraph.EdgeSystem.endpointWalkCoefficientSum; Komorebi.PaperIII.EdgeOperators.PacketGraph.EdgeSystem.transferOp_pow_apply_eq_sum_walks}

\begin{proof}
Take the $(j,i)$ entry in \cref{thm:word-to-walk}.  By
\cref{lem:word-endpoint}, a surviving word contributes
$c^\leftarrow(w)$ exactly when its endpoints are $i$ and $j$, and contributes
zero otherwise.  At $n=0$, the identity matrix gives the unique empty
diagonal walk.
\end{proof}

\begin{corollary}[Reachability obstruction at fixed length]
\label{cor:fixed-reachability}
If $\Walk_n(i,j)=\varnothing$, then $(T^n)_{ji}=0$.
\end{corollary}
\statusleanthree{Komorebi.PaperIII.EdgeOperators.PacketGraph.EdgeSystem.transferOp_pow_apply_eq_zero_of_no_walk}

\begin{proof}
This is immediate from the empty sum in the preceding corollary.
\end{proof}

\begin{corollary}[Orbit support obstruction]
\label{cor:orbit-reachability}
Let $S\subseteq V$, and let $x\in R^V$ be supported on $S$.  If $j$ is not
reachable by a directed path from any vertex of $S$, then
\[
  (T^nx)_j=0\qquad\text{for every }n\in\N.
\]
Hence the orbit under $T$ of $x$ cannot generate a vector with a nonzero
coordinate at $j$.
\end{corollary}
\statusleanthree{Komorebi.PaperIII.EdgeOperators.PacketGraph.EdgeSystem.Reachable; Komorebi.PaperIII.EdgeOperators.PacketGraph.EdgeSystem.transferOp_pow_mulVec_apply_eq_zero_of_unreachable_support}
\statusleanprior{Komorebi.PacketNaturalityCompletion.Reachability.algebraicOrbitHull_le_reachable_support; Komorebi.PacketNaturalityCompletion.Reachability.unreachable_nonzero_packet_obstruction}

\begin{proof}
Expand $(T^nx)_j=\sum_i(T^n)_{ji}x_i$.  Terms with $i\notin S$ vanish by
the support of $x$, and terms with $i\in S$ vanish by
\cref{cor:fixed-reachability}.  The final claim follows because the
$j$-coordinate also vanishes on every finite linear combination of orbit
vectors.
\end{proof}

\section{Canonical exact coordinate transport}
\label{sec:canonical-transport}

\subsection{The Paper III frame model}

Let every chart $a$ carry a linear frame equivalence
$F_a:M\xrightarrow{\sim}M$.  The formal Paper III module defines transport
from $a$ to $b$ by
\[
J_{b\leftarrow a}=F_b^{-1}F_a.
\]
This frame formula means \emph{synthesize in chart \(a\), then
analyze in chart \(b\)}.

\begin{theorem}[Exact frame transport is coherent]
\label{thm:exact-frame-transport}
For every $a,b,c$,
\[
J_{a\leftarrow a}=\id,\qquad
J_{a\leftarrow b}J_{b\leftarrow a}=\id,\qquad
J_{c\leftarrow b}J_{b\leftarrow a}=J_{c\leftarrow a},
\]
and therefore
\[
J_{a\leftarrow c}J_{c\leftarrow b}J_{b\leftarrow a}=\id.
\]
\end{theorem}
\statusleanthree{Komorebi.PaperIII.Transport.exactTransport_self; Komorebi.PaperIII.Transport.exactTransport_roundTrip; Komorebi.PaperIII.Transport.exactTransport_cocycle; Komorebi.PaperIII.Transport.exactTransport_triangleLoop}

\begin{proof}
Every identity is cancellation of adjacent $F_rF_r^{-1}$ factors.  The Lean
definitions use \leanname{LinearEquiv.trans}; the displayed formulas use the
equivalent right-to-left function-composition notation fixed in this paper.
\end{proof}

\begin{theorem}[Every finite canonical path telescopes]
\label{thm:exact-path-transport}
Let $p=(a_0,a_1,\ldots,a_n)$ be any finite vertex list and put
\[
J_p
=J_{a_n\leftarrow a_{n-1}}\cdots J_{a_1\leftarrow a_0},
\qquad
J_{(a_0)}=\id.
\]
Then
\[
J_p=J_{a_n\leftarrow a_0}.
\]
In particular, if $a_n=a_0$, then $J_p=\id$.
\end{theorem}
\statusleanthree{Komorebi.PaperIII.Transport.pathEndpoint; Komorebi.PaperIII.Transport.exactPathTransport; Komorebi.PaperIII.Transport.exactPathTransport_eq_endpoint; Komorebi.PaperIII.Transport.exactPathTransport_loop}

\begin{proof}
Induct on the list of successive vertices and apply the cocycle in
\cref{thm:exact-frame-transport} at each step.  If the endpoint is the start,
the endpoint transport is the identity.
\end{proof}

\subsection{Exact packet coordinates}

The graph model above resolves a fixed scalar matrix.  Coordinate transport
is a separate construction.  Let $X$ be an $R$-module, and let
\[
\cP^a=(P^a_i)_{i\in I_a}
\]
be a finite complete family of pairwise orthogonal linear idempotents:
\[
P^a_iP^a_j=\delta_{ij}P^a_i,
\qquad
\sum_{i\in I_a}P^a_i=I_X.
\]
Define the exact packet coordinate module
\[
\Coord(\cP^a)
=
\left\{(x_i)_{i\in I_a}:P^a_ix_i=x_i\right\}.
\]
Its analysis and synthesis maps are
\[
\cA_a(x)=(P^a_ix)_i,
\qquad
\cS_a((x_i)_i)=\sum_ix_i.
\]
Completeness and orthogonality give
\[
\cS_a\cA_a=\id_X,
\qquad
\cA_a\cS_a=\id_{\Coord(\cP^a)}.
\]

\begin{keyidea}[title={What a packet coordinate contains}]
A packet coordinate is a whole vector fixed by its projector. Analysis splits
an ambient vector into these components, and synthesis adds them back. The
transport between exact systems keeps the ambient vector and changes its
packet presentation.
\end{keyidea}

\begin{definition}[Canonical exact coordinate transport]
\label{def:canonical-J}
For two exact packet systems on the same ambient module, define
\[
  J^X_{b\leftarrow a}
  =
  \cA_b\cS_a:
  \Coord(\cP^a)\longrightarrow\Coord(\cP^b).
\]
The arrow points from the source coordinate system $a$ to the target
coordinate system $b$.
\end{definition}
\statusleanprior{Komorebi.PaperII.NestedRefinement.refineCoordinates; Komorebi.PaperII.NestedRefinement.coordinateRefinementEquiv}

\begin{theorem}[Identity, inverse, and cocycle for canonical transport]
\label{thm:J-cocycle}
For any exact packet systems $a,b,c$ on $X$,
\[
J^X_{a\leftarrow a}=\id,
\qquad
\left(J^X_{b\leftarrow a}\right)^{-1}=J^X_{a\leftarrow b},
\qquad
J^X_{c\leftarrow b}J^X_{b\leftarrow a}=J^X_{c\leftarrow a}.
\]
\end{theorem}
\statusleanprior{Komorebi.PaperII.NestedRefinement.refineCoordinates_self; Komorebi.PaperII.NestedRefinement.coarsenCoordinates_refineCoordinates; Komorebi.PaperII.NestedRefinement.refineCoordinates_coarsenCoordinates; Komorebi.PaperII.NestedRefinement.refineCoordinates_comp; Komorebi.PaperII.NestedRefinement.coordinateRefinementEquiv_trans}

\begin{proof}
Insert $\cS_b\cA_b=\id_X$:
\[
J^X_{c\leftarrow b}J^X_{b\leftarrow a}
=\cA_c\cS_b\cA_b\cS_a
=\cA_c\cS_a
=J^X_{c\leftarrow a}.
\]
The identity statement is $\cA_a\cS_a=\id$.  Setting $c=a$ in the cocycle
gives $J^X_{a\leftarrow b}J^X_{b\leftarrow a}=\id$; interchanging $a$ and $b$
gives the other composition, so the two transports are inverses.
\end{proof}

If $F:X\to Y$ is linear and $\cP^a,\cP^b$ are source packet systems while
$\cQ^r,\cQ^s$ are target packet systems, put
\[
\widehat F_{r\leftarrow a}
=
\cA^Y_rF\cS^X_a.
\]

\begin{theorem}[Canonical coordinate naturality]
\label{thm:canonical-naturality}
The exact coordinate representation preserves typed composition and satisfies
\[
J^Y_{s\leftarrow r}\widehat F_{r\leftarrow a}
=
\widehat F_{s\leftarrow b}J^X_{b\leftarrow a}.
\]
\end{theorem}
\statusleanprior{Komorebi.PaperII.NestedRefinement.coordinateOperator_comp; Komorebi.PaperII.NestedRefinement.coordinateOperator_naturality}

\begin{proof}
Both sides reduce to $\cA^Y_sF\cS^X_a$ after inserting the appropriate
synthesis--analysis identity.  The composition law is the same calculation
with the middle ambient module inserted.
\end{proof}

\begin{corollary}[Canonical triangle holonomy is trivial]
\label{cor:J-flat}
For every triple $a,b,c$,
\[
J^X_{a\leftarrow c}J^X_{c\leftarrow b}J^X_{b\leftarrow a}=\id.
\]
\end{corollary}
\statusleanprior{Komorebi.PaperII.NestedRefinement.refineCoordinates_comp; Komorebi.PaperII.NestedRefinement.refineCoordinates_self}

\begin{boundary}[Refinement is not reduction]
Although Paper II uses the name \emph{refineCoordinates}, $J$ is an
equivalence between two complete coordinate presentations.  It neither
truncates packets nor projects to a smaller state space.  Any lossy
restriction, interpolation, or learned identification is extra data and
cannot inherit \cref{thm:J-cocycle} without proof.
\end{boundary}

